% Re-authored after migration because the original notebook content was incomplete.

\chapter{積分の基礎}


\section{この章の目的}

微分が「ある点での変化の速さ」を扱うのに対し，積分は「小さな量を足し合わせて全体量を求める」ための道具です．
機械学習や信号処理では，
\begin{itemize}
\item 面積や総量の評価
\item 平均値や期待値の計算
\item 連続時間信号のエネルギー評価
\item 確率密度関数の正規化
\end{itemize}
などで積分が現れます．

本章の到達目標は次の通りです．
\begin{itemize}
\item 不定積分を「微分の逆演算」として理解できる
\item 定積分を「区間上の総和の極限」として理解できる
\item 広義積分が必要になる場面を説明できる
\item 二重積分の意味を面積・体積の観点から説明できる
\item 数値積分で近似的に積分値を求められる
\end{itemize}


\section{不定積分}

関数 $F(x)$ が関数 $f(x)$ の原始関数であるとは，
\[
F'(x) = f(x)
\]
が成り立つことを言います．

このとき，$f(x)$ の不定積分を
\[
\int f(x)\,dx = F(x) + C
\]
と書きます．
ここで $C$ は積分定数です．

積分定数が必要なのは，微分すると定数項が消えるからです．
例えば
\[
\frac{d}{dx}(x^3 + 7) = 3x^2, \quad
\frac{d}{dx}(x^3 - 2) = 3x^2
\]
なので，$3x^2$ の原始関数は 1 つではありません．

\subsection{基本公式}

微分の公式を逆向きに使うと，以下の不定積分公式が得られます．
\[
\int x^n\,dx = \frac{x^{n+1}}{n+1} + C \quad (n \neq -1)
\]
\[
\int \cos x\,dx = \sin x + C
\]
\[
\int \sin x\,dx = -\cos x + C
\]
\[
\int \frac{1}{x}\,dx = \log |x| + C
\]

ここで最後の式だけ形が違うのは，$x^{-1}$ に対して上のべき乗則をそのまま適用できないからです．

\subsection{線形性}

不定積分も線形性を持ちます．定数 $a, b$ と関数 $f, g$ に対して
\[
\int \left(af(x) + bg(x)\right)\,dx
= a\int f(x)\,dx + b\int g(x)\,dx
\]
が成り立ちます．

複雑な積分を基本公式の組み合わせに分解する際には，まず線形性で項をばらすのが基本です．


\section{定積分}

\subsection{面積としての理解}

区間 $[a,b]$ 上で非負の関数 $f(x)$ を考えるとき，
\[
\int_a^b f(x)\,dx
\]
は，グラフと $x$ 軸に挟まれた領域の面積として理解できます．

ただし積分は単なる「面積」より少し一般的で，関数が負になる部分は負の寄与として数えます．
したがって定積分は，\textbf{符号付き面積}と考えると正確です．

\subsection{リーマン和}

定積分の本質は，「区間を細かく分けて長方形の面積を足し合わせ，分割を無限に細かくした極限を取る」ことです．
区間 $[a,b]$ を $N$ 分割し，幅を
\[
\Delta x = \frac{b-a}{N}
\]
とします．
各小区間の代表点を $x_k$ とすると，リーマン和
\[
\sum_{k=0}^{N-1} f(x_k)\Delta x
\]
が定積分の近似になり，
\[
\int_a^b f(x)\,dx
= \lim_{N\to\infty}\sum_{k=0}^{N-1} f(x_k)\Delta x
\]
で定義されます．

この見方は数値積分や確率論で極めて重要です．
期待値や連続時間信号のエネルギーも，本質的には「無数の小さな寄与の総和」です．

\subsection{微積分の基本定理}

定積分と不定積分をつなぐ最重要結果が\textbf{微積分の基本定理}です．
$F'(x)=f(x)$ を満たす原始関数 $F$ があるとき，
\[
\int_a^b f(x)\,dx = F(b) - F(a)
\]
が成り立ちます．

この定理により，面積の極限として定義された量を，原始関数の差として効率よく計算できます．

例えば
\[
\int_0^1 x^2\,dx
= \left[\frac{x^3}{3}\right]_0^1
= \frac{1}{3}
\]
です．


\section{置換積分と部分積分の考え方}

本章では詳細な技巧を網羅しませんが，よく使う発想だけは押さえておきます．

\subsection{置換積分}

合成関数の積分では，新しい変数に置き換えて簡単な形に直すことがあります．
例えば
\[
\int 2x \cos(x^2)\,dx
\]
では $u = x^2$ と置くと $du = 2x\,dx$ なので，
\[
\int \cos u\,du = \sin u + C = \sin(x^2) + C
\]
と計算できます．

\subsection{部分積分}

積の微分公式
\[
(fg)' = f'g + fg'
\]
を変形すると
\[
\int f(x)g'(x)\,dx
= f(x)g(x) - \int f'(x)g(x)\,dx
\]
が得られます．

これは「微分すると簡単になる部分」と「積分しやすい部分」を分ける操作です．
例えば多項式と指数関数，$x\sin x$ のような積で有効です．


\section{広義積分}

普通の定積分は有限区間上の有界な関数を想定していますが，実際には
\begin{itemize}
\item 積分区間が無限に広い
\item 被積分関数が途中で発散する
\end{itemize}
場合もあります．
このような積分を\textbf{広義積分}と呼びます．

例えば
\[
\int_1^\infty \frac{1}{x^2}\,dx
\]
は上端が無限大なので，極限で
\[
\int_1^\infty \frac{1}{x^2}\,dx
= \lim_{R\to\infty}\int_1^R \frac{1}{x^2}\,dx
= \lim_{R\to\infty}\left[-\frac{1}{x}\right]_1^R
= 1
\]
と定義します．

一方，
\[
\int_1^\infty \frac{1}{x}\,dx
\]
は
\[
\lim_{R\to\infty}\int_1^R \frac{1}{x}\,dx
= \lim_{R\to\infty}\log R
\]
となり発散します．

「無限区間だから必ず発散する」わけではなく，\textbf{どれくらい速く減衰するか}が重要です．


\section{多重積分}

2 変数関数 $f(x,y)$ の積分では，平面上の微小面積要素を足し合わせます．
長方形領域 $a \leq x \leq b,\ c \leq y \leq d$ 上では，
\[
\iint f(x,y)\,dx\,dy
\]
を考えます．

もし $f(x,y)$ がその点での高さを表すなら，二重積分は曲面の下の体積に対応します．

計算上は，1 変数ずつ順に積分することが多く，
\[
\iint_R f(x,y)\,dx\,dy
= \int_a^b \left(\int_c^d f(x,y)\,dy\right)dx
\]
のように扱います．

例えば単位正方形上の $f(x,y)=x+y$ なら
\[
\int_0^1 \int_0^1 (x+y)\,dy\,dx
= \int_0^1 \left(x + \frac{1}{2}\right)dx
= 1
\]
です．

画像処理では 2 次元信号，確率論では同時密度関数，機械学習では複数変数にまたがる平均化で多重積分が現れます．


\section{Python で数値積分を確かめる}

解析的に原始関数が求まらない場合や，離散データしかない場合には数値積分を使います．
以下は台形公式で積分値を近似する例です．

\begin{lstlisting}[style=codecell,language=Python]
from pathlib import Path

import matplotlib.pyplot as plt
import numpy as np

output_dir = Path("outputs/ch05")
output_dir.mkdir(parents=True, exist_ok=True)

x = np.linspace(0.0, 1.0, 1001)
y = x ** 2

approx = np.trapezoid(y, x)
print(f"approx integral = {approx:.6f}")

plt.figure(figsize=(4, 3))
plt.plot(x, y, label="y = x^2")
plt.fill_between(x, y, alpha=0.3)
plt.xlabel("x")
plt.ylabel("y")
plt.tight_layout()
plt.savefig(output_dir / "integral_x_squared.png", dpi=150)
plt.close()
\end{lstlisting}

この例では，$\int_0^1 x^2\,dx = 1/3$ に近い値が得られるはずです．
分割数を増やすと近似誤差が小さくなることも確認してください．


\section{この章で押さえるべき点}

\begin{itemize}
\item 不定積分は原始関数を求める操作であり，積分定数が必要である
\item 定積分は区間上の総和の極限であり，符号付き面積として理解できる
\item 微積分の基本定理により，定積分は原始関数の差で計算できる
\item 広義積分では無限区間や特異点を極限で扱う
\item 多重積分は高次元での総和であり，画像・確率・信号処理で広く使う
\end{itemize}

次の線形代数・確率統計の章では，積分で定義される平均値や期待値が頻出します．
そのとき，積分を「連続量の足し合わせ」として捉えられるかが理解の分かれ目です．
